% ============================================================
% 标题 + 摘要 + 关键词
% ------------------------------------------------------------
% 说明：全国大学生数学建模竞赛的封面由组委会统一提供（含题目、
% 参赛队号、密封线等），正文一般从「摘要」开始单独编排。
% 如需要自制封面，可在此处插入 titlepage，比赛提交时请按实际要求删除/替换。
% ------------------------------------------------------------

\begin{center}
  {\zihao{3}\bfseries [论文题目]}
\end{center}

\vspace{1em}
\begin{center}
  {\zihao{-2}\bfseries 摘\quad 要}
\end{center}

\noindent
% ------------------------------------------------------------
% 摘要写作要点（参考样例论文 C023）：
% 1. 一句话点出问题背景与研究对象；
% 2. 概述总体建模思路（用了什么方法框架）；
% 3. 分问题逐段说明「做了什么 —— 用什么方法 —— 得到什么结果（含关键数字）」；
% 4. 末段点出创新点 / 应用价值。
% ------------------------------------------------------------
【占位：一段总起，交代问题背景、研究对象与总体建模框架。】

\textbf{针对问题一，}【占位：方法 + 关键结果】。

\textbf{针对问题二，}【占位：方法 + 关键结果】。

\textbf{针对问题三，}【占位：方法 + 关键结果】。

\textbf{针对问题四，}【占位：方法 + 关键结果】。

【占位：总结全文创新点，如模型融合思路、算法设计、可视化/工程实现等。】

\vspace{0.6em}
\noindent\textbf{关键词：} [关键词1]\quad [关键词2]\quad [关键词3]\quad [关键词4]\quad [关键词5]

\newpage
